% =====================================================================
% Management Science (INFORMS) – Research Proposal LaTeX Template
% ---------------------------------------------------------------------
% Notes:
% • This is a proposal-style manuscript aimed at Management Science.
% • It follows common INFORMS/author–year conventions via natbib.
% • Replace placeholders (ALL-CAPS) with your content.
% • Swap bibliography style to an INFORMS .bst if required by your target track
%   (e.g., informs2014.bst) by changing \bibliographystyle below.
% =====================================================================

\documentclass[12pt]{article}

\usepackage{float}
% ----- Encoding, fonts, microtype
\usepackage[T1]{fontenc}
\usepackage[utf8]{inputenc}
\usepackage{lmodern}
\usepackage{textalpha}

\usepackage{microtype}

% ----- Layout & spacing
\usepackage[margin=1in]{geometry}
\usepackage{setspace}
\onehalfspacing   % or \doublespacing

% ----- Math & symbols
\usepackage{amsmath,amssymb}
\usepackage{textcomp}

% ----- PGFPlots & TikZ for figures
\usepackage{tikz}
\usepackage{pgfplots}
\pgfplotsset{compat=1.18}
\usetikzlibrary{arrows.meta, positioning}
\usepackage{enumitem} 
% ----- Tables & figures
\usepackage{booktabs}
\usepackage{threeparttable}
\usepackage{booktabs}      % for \toprule, \midrule, \bottomrule
    % for numeric alignment (optional)

\usepackage{siunitx}
\usepackage{array,tabularx}
\newcolumntype{Y}{>{\centering\arraybackslash}X}
\usepackage{graphicx}
\usepackage{caption}
\usepackage{subcaption}
\usepackage{adjustbox}
\pdfminorversion=7
\usepackage{media9}

% ----- Lists & colors & hyperlinks
\usepackage{xcolor}
\usepackage{enumitem}
\setlist[itemize]{leftmargin=2em,label=\textbullet}
\setlist[enumerate]{leftmargin=2em}
\usepackage[hidelinks]{hyperref}


% ----- Quotes & hyphenation
\usepackage{csquotes}
\usepackage{hyphenat}

% ----- Handy unicode chars
\usepackage{newunicodechar}
\DeclareUnicodeCharacter{2265}{\geq}
\DeclareUnicodeCharacter{00A5}{\textyen} % <-- fixes ¥
\newunicodechar{≈}{\approx}
\newunicodechar{−}{\textminus}

% ===== Bibliography =====
\usepackage[
    style=vancouver,
    backend=biber,
    sorting=none,      % citations appear in order of appearance
    doi=true,
    url=false,         % Vancouver style: no URLs for journal articles
    natbib=true
]{biblatex}

\addbibresource{reference.bib}

% Make citations SUPERSCRIPT like AHEHP wants
\usepackage{csquotes}
\renewcommand{\cite}{\supercite}



% ---------- Begin Document ----------

\begin{document}


\title{The Shadow Price of Vaccination Delay: Evidence From a Case Study from Wuhan, China}

\author{Blinded for peer review}
\maketitle

% ----- Abstract -----
\begin{abstract}
Vaccination timing critically determines the duration of biological exposure to infectious 
risk, yet many individuals delay immunisation even when vaccines are free and accessible. 
While traditional models focus on binary uptake, this study quantifies the behavioral cost 
of waiting as an intertemporal health investment. Using a large-scale discrete choice 
experiment (N=1,027) in Wuhan, China, we recover structural time-preference parameters 
characterizing sensitivity to delayed biological protection.

Results reveal pronounced present-oriented preferences: respondents strongly discount 
delayed vaccination, with hyperbolic parameter $\kappa = 0.225$ and exponential 
$\delta = 0.160$. These parameters translate into a marginal willingness-to-accept (MWTA) 
of 45--56 RMB per month of waiting. Delay aversion is systematically patterned by 
biological and social vulnerability—individuals with chronic conditions, lower 
institutional trust, and irregular preventive-health routines exhibit substantially 
steeper discounting and require higher compensation to tolerate delay.

By linking behavioral parameters to biological risk and uncertainty, we show that 
vaccination delays impose regressive behavioral burdens. We integrate these parameters 
into a stylized epidemiological framework to simulate how heterogeneous delay aversion 
impacts outbreak trajectories. The findings provide actionable inputs for designing 
time-consistent public health strategies that minimize the behavioral cost of waiting and 
improve the efficiency and equity of preventive-health delivery.

\textbf{Keywords:} Vaccination timing; intertemporal preferences; biological vulnerability; 
discrete choice experiment; SEIR modeling; willingness-to-accept; delay aversion; China.
\end{abstract}




% ----- Proposal Overview -----
\section{Introduction}
\doublespacing

Timely access to biological protection is a central determinant of human exposure to 
infectious disease and the preservation of health capital \citep{grossman1972concept, 
becker1964human}. While vaccination coverage is often treated as a binary outcome---vaccinated 
or not---the timing of vaccination critically shapes the duration for which individuals 
remain unprotected and biologically vulnerable. Delays in vaccination therefore represent a 
protracted interval of continued exposure, during which individuals face accumulating 
infection risk and potential health deterioration. Understanding how humans evaluate such 
delays is essential for explaining patterns of preventive health behavior under infectious 
risk.

From an intertemporal perspective, vaccination timing represents a health investment decision 
in which immediate costs---such as effort, inconvenience, or perceived side effects---are 
traded off against delayed and uncertain biological benefits \citep{cropper1977health}. 
When protection is postponed, individuals must decide whether the future reduction in 
infection risk sufficiently compensates for remaining exposed in the present. These 
trade-offs are not uniform. Perceived biological vulnerability, institutional trust, and 
stable preventive-health routines shape how future protection is perceived and discounted, 
giving rise to systematic heterogeneity in tolerance for delayed protection.

Despite the clinical and epidemiological importance of timing, much of the behavioral 
literature treats waiting time as an administrative friction rather than as a biologically 
meaningful interval of exposure. This study addresses this gap by quantifying how individuals 
value timely protection as an intertemporal health trade-off. Using a discrete choice 
experiment (DCE) administered to 1,027 adults in Wuhan, China, we recover structural 
time-preference parameters characterizing sensitivity to delayed biological protection. 
By embedding waiting time directly into vaccine attributes, we estimate both the hyperbolic 
($\kappa$) and exponential ($\delta$) discount parameters and convert these into a marginal 
willingness-to-accept (MWTA) for delay.

Wuhan provides a particularly informative setting for studying these processes. As a city 
with high population density and direct experience with rapid epidemic surges, Wuhan 
constitutes a high-salience environment where the biological consequences of delayed 
protection are readily understood by respondents. This salience sharpens the behavioral 
signal linking perceived vulnerability, trust, and tolerance for waiting. While absolute 
valuations of delay may vary across contexts, the underlying mechanisms identified 
here---risk-weighted time preferences, present-oriented valuation of protection, and 
heterogeneity by biological and social vulnerability---reflect general features of human 
health behavior rather than idiosyncratic local conditions.

We make three contributions. First, we provide a robust quantification of intertemporal 
preferences for vaccination, recovering pronounced short-run impatience that aligns with 
present-biased health investment models. Second, we demonstrate that sensitivity to delay is 
systematically patterned by biological characteristics, indicating that time preferences 
are shaped by perceived vulnerability and uncertainty. Third, we integrate these behavioral 
parameters into a stylized epidemiological framework to simulate the population-level 
impact of heterogeneous delay aversion. Together, these findings position vaccination delay 
as a biologically consequential dimension of preventive behavior and underscore the 
importance of accounting for intertemporal preferences when designing public health 
strategies.


\section{Theoretical Framework}

Vaccination timing can be conceptualized as an investment in health capital, in which 
individuals weigh the immediate utility costs of uptake against the discounted preservation 
of future health stocks \citep{grossman1972concept}. In this framework, postponing 
protection extends the period during which individuals remain biologically exposed to 
infectious risk. The discount parameters therefore reflect valuations shaped by individuals' 
biological vulnerability and institutional environment.

We characterize intertemporal preferences using two standard specifications. Under an 
exponential specification, the utility of delayed protection at time $t$ is given by 
$U_t = U_0 \cdot \exp(-\delta t)$, where $\delta$ reflects a constant rate of time preference. 
To account for the possibility of pronounced short-term impatience, we also employ a 
quasi-hyperbolic specification, $U_t = \frac{U_0}{1 + \kappa t}$, where $\kappa$ captures 
present-oriented discounting. In our context, $\kappa$ and $\delta$ represent not merely 
abstract time preferences, but valuations of remaining unprotected while exposure risk 
persists.

Discrete choice experiments (DCEs) provide a structured approach to quantifying these 
intertemporal valuations. The latent utility associated with a vaccine alternative with 
delay $t$ and monetary incentive $C$ can be expressed as:
\[
V(t, C) = \beta_{delay} \cdot t + \beta_{cash} \cdot C + \dots
\]
The ratio $-\beta_{delay} / \beta_{cash}$ yields the marginal willingness-to-accept (MWTA) 
for one additional unit of delay. This MWTA captures the "shadow price" of timely protection. 
Importantly, sensitivity to cost ($\beta_{cash}$) may vary across the population due to 
varying budget constraints and income levels, making MWTA a relative measure of the 
behavioral burden of waiting.

Uncertainty further amplifies the behavioral cost of delay. When the future state of the 
epidemic or vaccine effectiveness is ambiguous, individuals tend to discount delayed 
outcomes more heavily \citep{Attema2013}. Lower institutional trust may heighten this 
perceived uncertainty, increasing impatience toward delayed protection \parencite{Liu2021}. 
Anticipated heterogeneity across demographic and health-related groups reflects differences in 
vulnerability and trust, allowing identification of populations for whom vaccine delays are 
especially costly. 


\section{Method}

\subsection{Data and Experimental Design}

Data were obtained from a discrete choice experiment (DCE) conducted among 1,027 adults in 
Wuhan, China, between March and May 2025. 
 Wuhan is a theoretically relevant setting because 
its early COVID-19 experience, intensive public-health response, and rapid vaccine rollout 
made issues of risk, institutional trust, and government capacity especially salient—key 
determinants of delay sensitivity and intertemporal vaccination behaviour. At the same time, 
these features limit broader generalizability: regions with lower institutional visibility, 
different pandemic trajectories, or weaker health systems may exhibit different delay–cash 
trade-offs. Accordingly, the estimates should be interpreted within the Wuhan context.

Respondents were asked to consider a future “COVID-like disease” (CLD), defined as a 
hypothetical respiratory infection with severity and transmission similar to COVID-19. This 
neutral framing enabled us to elicit vaccination preferences without anchoring responses to 
any specific COVID-19 wave or variant. Stated-preference methods provide a rigorous way to quantify delay–cash trade-offs that 
cannot be directly observed in revealed-preference data or ethically induced through 
experimental variation in real vaccination programmes. The DCE was designed and implemented 
in accordance with ISPOR good-practice guidelines for attribute development, experimental 
design, pilot testing, and analysis, ensuring the validity and policy interpretability of  
the estimated behavioural parameters.



Each participant completed six choice tasks involving two hypothetical vaccines for the CLD 
and an opt-out option. Vaccine profiles varied in delay, efficacy, side effects, origin, and 
monetary incentives, enabling estimation of cash–delay trade-offs and associated behavioural 
parameters. Appendix~A provides a comparison of the analytic sample with 2020 Wuhan Census 
benchmarks.

The full experimental protocol, sampling procedures, and ethical safeguards were implemented
following standard guidelines for stated-preference research in health economics and were
approved by the relevant institutional review boards.
 Although the survey instrument is shared with prior work, the present 
study pursues distinct aims: estimating marginal willingness-to-accept (MWTA) for waiting, 
analysing behavioural impatience, and conducting new policy simulations based on revised 
economic specifications.

\subsection{Development of Attributes and Levels}

Attribute development followed established guidance for stated-preference research in health 
economics and aimed to capture the core trade-offs individuals face in real vaccination 
decisions. Candidate attributes were identified through a review of vaccination behavior 
studies, policy reports, and behavioral-economic frameworks emphasising delay, risk, and 
incentives. Consultations with public-health practitioners in China and a 60-respondent pilot 
survey refined attribute wording, ensured cognitive clarity, and confirmed that the attributes 
reflected realistic decision environments.

The final attributes included vaccine delays (0, 1, 3, 6 months), vaccine efficacy (50\%, 70\%, 
95\%), side-effect severity (mild, moderate, severe), monetary incentives (0, 50, 200, 800 RMB), 
and vaccine origin (domestic, imported). Vaccine-delay levels were chosen to reflect the 
typical postponements generated by supply fluctuations, phased eligibility, and appointment 
bottlenecks in large-scale rollouts. Pretest findings indicated that most respondents viewed 
delays of 1--3 months as realistic and acceptable, confirming that monthly increments 
represent a salient and credible unit over which individuals evaluate delayed protection.

Monthly increments also align with the practical realities of vaccination planning, where 
delays of one to several months—not days—are the operationally relevant horizon over which 
programmes adjust supply distribution, staffing, and booking systems. Using monthly delay 
levels therefore enables estimation of the behaviourally meaningful compensation required to 
tolerate the types of postponements that actually occur in real programmes.

Methodologically, this implies that the DCE is not intended to recover very short-run 
discounting curvature or identify structural $\beta$–$\delta$ parameters, which require day-level 
variation. Instead, the experiment yields reduced-form valuations of vaccine delay over a 
mid-range horizon, providing estimates that align directly with the operational decisions 
policymakers face when managing congestion and scheduling in vaccination programmes.

Financial-incentive levels were informed by magnitudes used in recent field experiments and 
policy pilots in China and neighbouring East Asian jurisdictions, ensuring that respondents 
evaluated realistic and policy-relevant amounts. Each choice task included an opt-out option, 
which enhances realism and aligns with good practice in DCE design.

The pilot study confirmed that respondents understood the tasks, that no attribute dominated 
choices, and that minor wording revisions improved clarity and reduced cognitive burden 
prior to fielding the main survey.

\begin{figure}[h]
    \centering
    \includegraphics[width=0.95\textwidth,height=0.8\textheight,keepaspectratio]{Code_Generated_Image (2).png}
    \caption{
        Development of attributes and levels for the discrete choice experiment.
        The diagram summarises the four-step process used to identify candidate
        attributes, refine wording through pilot testing, justify level ranges,
        and finalise the attribute set following ISPOR good-practice guidelines.
    }
    \label{fig:attribute_development}
\end{figure}

\clearpage

\medskip

\begin{table}[htbp]
\centering
\captionsetup{font=small}
\caption{Experimental Attributes, Levels, and Coding Used in the Discrete Choice Experiment}
\small
\begin{tabularx}{0.9\textwidth}{@{} l Y Y Y @{}}
\toprule
\textbf{Attribute} & \textbf{Levels} & \textbf{Coding} & \textbf{Description} \\
\midrule
Vaccine Delays &
0, 1, 3, 6 months (Walk-in to 6 mo) &
0, 1, 3, 6 &
Waiting time before the vaccine becomes available. \\

Vaccine Origin &
Domestic; Imported &
0 = Domestic, 1 = Imported &
Indicates whether the vaccine was produced locally or overseas. \\

Vaccine Efficacy &
50\%, 70\%, 95\% &
0.50, 0.70, 0.95 &
Probability of protection against infection. \\

Side Effects &
Mild; Moderate; Severe &
0, 1, 2 &
Ordered severity of expected vaccine side effects. Higher values indicate more severe reactions. \\

Monetary Incentives (RMB) &
50, 200, 800 &
50, 200, 800 &
Monetary compensation offered upon receiving the vaccine. \\
\bottomrule
\end{tabularx}

\vspace{0.25em}
\captionsetup{font=footnotesize}
\caption*{\textit{Notes:} Attribute codings reflect the numerical values entered into the econometric model. Side-effect severity is treated as an ordinal variable.}
\end{table}


\clearpage



\clearpage

% =======================
% 2.2 Modeling and Estimation
% =======================
\subsection{Modeling and Estimation}

\subsubsection{Model Framework}

We model the observed choice outcome as the dependent variable to quantify the cash--delay
trade-offs underlying vaccination decisions. For each respondent $i$ ($i = 1,\ldots,N$),
choice alternative $j$ ($j \in \{A,B,C\}$), and choice task $t$ ($t = 1,\ldots,T_i$), the
indicator $Y_{ijt}$ is defined as
\[
Y_{ijt} =
\begin{cases}
1, & \text{if respondent $i$ selects vaccine $j$ in choice task $t$},\\[4pt]
0, & \text{otherwise.}
\end{cases}
\]

The probability of choosing alternative $j$ is assumed to depend on a latent utility
function $U_{ijt}$, which consists of a systematic component $V_{ijt}$---a linear
combination of vaccine attributes---and an unobserved stochastic term $\varepsilon_{ijt}$:
\[
U_{ijt} = V_{ijt} + \varepsilon_{ijt}, \qquad
V_{ijt} = \alpha_j
+ \beta_{\text{delay}} \,\text{Delay}_{ijt}
+ \beta_{\text{cash}}  \,\text{Cash}_{ijt}
+ \beta_{\text{eff}}   \,\text{Eff}_{ijt}
+ \beta_{\text{se}}    \,\text{SE}_{ijt}
+ \beta_{\text{orig}}  \,\text{Orig}_{ijt},
\]
where $\alpha_C \equiv 0$ and $\alpha_A, \alpha_B$ are alternative-specific constants (ASCs).

\noindent \textbf{Attribute definitions:}
\begin{itemize}[leftmargin=1.5em]
  \item \textit{Delay:} waiting time before vaccination (months).
  \item \textit{Cash:} monetary incentive (RMB), coded per~¥100.
  \item \textit{Eff:} vaccine efficacy level (0.50, 0.70, 0.95).
  \item \textit{SE:} side-effect severity, coded as an ordered variable
        ($0 = \text{Mild}$, $1 = \text{Moderate}$, $2 = \text{Severe}$), with higher values
        indicating more severe side effects.
  \item \textit{Orig:} vaccine origin indicator ($1 = \text{Imported}$; $0 = \text{Domestic}$).
\end{itemize}

Assuming that the unobserved error terms $\varepsilon_{ijt}$ are independently and
identically distributed with a Type~I extreme-value distribution, the probability that
respondent $i$ chooses alternative $j$ in task $t$ follows the multinomial logit (MNL) form:
\[
P_{ijt} = \Pr(Y_{it} = j)
= \frac{\exp(V_{ijt})}{\sum_{k \in \{A,B,C\}} \exp(V_{ikt})}.
\]

The log-likelihood function to be maximised is
\[
L(\theta)
= \sum_i \sum_t \sum_j \mathbf{1}\{Y_{it}=j\} \, \log P_{ijt}(\theta),
\]
where $\mathbf{1}\{\cdot\}$ equals one if alternative $j$ is chosen in task $t$ and zero otherwise.

% -------------------------------------------------
\subsubsection{Estimation and Coefficients}

The parameter vector $\hat{\beta}$ was estimated via maximum likelihood under the MNL
specification. Estimated coefficients represent marginal utilities associated with each
vaccine attribute: a negative $\hat{\beta}_{\text{delay}}$ reflects time aversion, while a
positive $\hat{\beta}_{\text{cash}}$ indicates that monetary incentives increase vaccination
uptake. All models were estimated using the \texttt{mlogit} package in \texttt{R} (version
4.3), with standard errors clustered at the individual level to account for repeated choices
per respondent.

% -------------------------------------------------
\subsubsection{Marginal Willingness to Accept (MWTA)}

To facilitate policy interpretation, coefficients were converted into monetary equivalents.
The marginal willingness to accept (MWTA) for one additional month of delay is defined as
\[
\text{MWTA}_{\text{delay}} = -\frac{\hat{\beta}_{\text{delay}}}{\hat{\beta}_{\text{cash}}},
\]
representing the monthly compensation required to offset the disutility of waiting. Higher
MWTA values therefore indicate stronger impatience toward delayed vaccination.

MWTA is a natural and policy-relevant metric in this setting because both delay and monetary
incentives are concrete, easily understood attributes that respondents can evaluate
credibly. Methodological work on health-related discrete choice experiments shows that MWTA
values derived from such designs closely track real-world trade-offs when attributes are
objective and quantifiable \citep{Johnson2013,Veldwijk2019}. In our context, MWTA provides a
direct behavioural valuation of waiting time and a transparent basis for setting compensation
levels or comparing investments in reducing delays versus increasing incentives.

% -------------------------------------------------
\subsubsection{Counterfactual Uptake Under Policy Bundles}

To translate these behavioural parameters into policy-relevant quantities, we simulate
predicted vaccine uptake under counterfactual policy bundles that vary waiting time
(0--6~months) and monetary incentives (0--800~RMB), holding efficacy at 70\%, side effects
at mild, and origin as domestic. For each bundle, the predicted uptake probability is
computed as
\[
\hat{P}_j(\text{bundle}) = 
\frac{\exp(V_j(\hat{\beta}))}
     {\sum_{k \in \{A,B,C\}} \exp(V_k(\hat{\beta}))},
\]
where $V_j(\hat{\beta})$ denotes the deterministic utility of vaccine $j$ evaluated at the
estimated coefficients. The resulting probabilities, summarised in Table~5, quantify how
monetary incentives compensate for longer vaccination delays and provide inputs for the
policy simulations reported in Section~5.

% -------------------------------------------------
\subsection{Validity Checks}

The credibility of the stated-preference responses was assessed using internal, external,
and construct validity checks, with detailed diagnostics reported in Appendix~B.

Internal validity (Appendix~B.1) is supported by high dominance-test pass rates, strong
repeated-choice consistency, and coefficient stability in sensitivity analyses excluding
respondents who may have engaged in attribute non-attendance. These indicators suggest that
respondents attended to the attributes and made coherent, compensatory trade-offs across
choice tasks.

External validity (Appendix~B.2) was assessed by comparing the model’s predicted incentive
elasticities with real-world evidence from vaccination incentive programmes in China
\parencite{Hu2022,Li2023VaccIncentives}, Hong Kong \parencite{Yip2022HKincentives}, South
Korea \parencite{Kim2022SKincentive}, Sweden \parencite{CamposMercade2021}, and the United
States \parencite{Dave2021USincentives}. Field evaluations in Henan, Guangdong, Shanghai,
and Hong Kong document increases of 3--15 percentage points for incentives valued at
10--450~RMB \parencite{Hu2022,Li2023VaccIncentives,Yip2022HKincentives}, with comparable
effect sizes reported in South Korea and Europe \parencite{Kim2022SKincentive,
CamposMercade2021}. Despite differences in programme design and baseline uptake, these
empirical magnitudes fall within the range predicted by our DCE, supporting the external
validity of the stated-preference estimates.

Construct validity (Appendix~B.3) is reflected in coefficient signs, magnitudes, and subgroup
gradients that align with theoretical expectations and prior evidence on vaccination
behaviour, time preferences, and institutional trust. Finally, methodological considerations
(Appendix~B.4) underscore that multi-month delays cannot be experimentally manipulated or
observed cleanly in revealed-preference data, making a discrete choice experiment the most
feasible and ethically appropriate approach for estimating delay--cash trade-offs.

Taken together, these checks support the reliability, interpretability, and applied policy
relevance of the behavioural parameters used in the subsequent analyses and simulations.






\clearpage

\section{Results}

Table~2 summarizes the estimated effects the attributes on choice probabilities under the Multinomial Logit (MNL) and Mixed Logit (MXL) specifications. Standardized coefficients are reported for comparability across attributes, since variables measured in different units (e.g., months, RMB, percentage points) were scaled to ensure coefficient magnitudes reflect relative importance rather than unit size. In both model, coefficients on vaccine dealy are negative and highly significant, indicating that respondents strongly dislike waiting for vaccination.
 The corresponding positive and significant coefficients on the monetary incentives in both models confirm that higher monetary rewards substantially increase the likelihood of vaccine uptake.Vaccine efficacy is also positively valued, consistent with the expectation that greater protection enhances perceived utility. In contrast, side effects show a small and statistically weak influence in the MNL model, suggesting that concerns about adverse reactions play a limited role once other attributes are controlled. Respondents exhibit a strong preference for imported vaccines in the MNL estimation, although this effect weakens under the MXL estimation. Overall, the results showing that reducing vaccine delay and providing monetary incentives are behaviorally effective levers to accelerate vaccination uptake.

In this study, the multinomial logit (MNL) model is used as the primary specification. While 
mixed logit (MXL) models are frequently employed to account for unobserved preference 
heterogeneity in DCEs, several considerations make the MNL particularly appropriate for our 
setting. First, the MNL offers a transparent interpretation of the cash and delay coefficients, 
which is essential for deriving marginal willingness-to-accept (MWTA) measures and conducting 
policy simulations. Second, each respondent completed only six choice tasks, limiting the 
empirical stability of highly parameterised random-coefficients models; preliminary MXL 
estimations produced imprecise or unstable variance components for multiple attributes, a 
well-documented challenge in short DCE panels. Third, as shown in Appendix~B 3.3, the MXL results 
yield coefficient patterns and MWTA estimates that closely mirror those from the MNL model, 
indicating that our substantive behavioural conclusions are robust to more flexible 
preference structures. Together, these considerations support the use of the MNL model as the 
main analytic framework.


\clearpage
\begin{table}[htbp]
\centering
\begin{threeparttable}
\caption{Main model estimates (standardised attributes): multinomial logit (MNL) and mixed logit (MXL)}
\label{tab:main_models}
\setlength{\tabcolsep}{3.8pt}
\footnotesize
\begin{tabular}{lccccc}
\toprule
& \multicolumn{2}{c}{MNL} & \multicolumn{3}{c}{MXL} \\
\cmidrule(lr){2-3}\cmidrule(lr){4-6}
Parameter
& Coef.\ (SE) & 95\% CI
& Mean (SE) & 95\% CI & SD (SE) \\
\midrule
ASC (B vs A)           & $-0.108^{***}$ $(0.032)$ & $[-0.170,\,-0.045]$ & ---                 & ---                 & ---           \\
ASC (C vs A)           & $-0.326^{***}$ $(0.062)$ & $[-0.448,\,-0.203]$ & ---                 & ---                 & ---           \\
Vaccine delays (std)    & $-0.125^{***}$ $(0.027)$ & $[-0.177,\,-0.072]$ & $-0.353^{***}$ $(0.027)$ & $[-0.405,\,-0.301]$ & $1.303$ $(0.073)$ \\
Monetary incentives (std)  & $ 0.313^{***}$ $(0.018)$ & $[ 0.278,\, 0.347]$ & $ 0.431^{***}$ $(0.025)$ & $[ 0.383,\, 0.479]$ & ---           \\
Vaccine efficacy (std) & $ 0.163^{***}$ $(0.025)$ & $[ 0.115,\, 0.212]$ & $ 0.218^{***}$ $(0.041)$ & $[ 0.138,\, 0.297]$ & ---           \\
Side effects (std)     & $-0.082^{*}$   $(0.033)$ & $[-0.147,\,-0.017]$ & $ 0.008$        $(0.031)$ & $[-0.052,\, 0.068]$ & ---           \\
Imported origin        & $ 0.239^{***}$ $(0.048)$ & $[ 0.145,\, 0.333]$ & $ 0.093^{*}$   $(0.047)$ & $[ 0.001,\, 0.185]$ & ---           \\
ASC (opt-out)          & ---                     & ---                 & $-0.287^{*}$   $(0.118)$ & $[-0.518,\,-0.056]$ & ---           \\
\bottomrule
\end{tabular}
\begin{tablenotes}[flushleft]
\footnotesize
\item Notes: Standardised coefficients are reported for comparability across attributes; variables
measured in different units (months, RMB, percentage points) were scaled prior to estimation.
Robust standard errors (SE) are clustered at the individual level. Confidence intervals (CI) are 
95\% Wald intervals. MXL estimates allow a random coefficient on vaccine delay; other 
coefficients are fixed. Sample size: 1{,}027 respondents (6 choice tasks per respondent, 3 
alternatives per task). 
\item[] $^{***} p<0.01$; $^{**} p<0.05$; $^{*} p<0.10$.
\end{tablenotes}
\end{threeparttable}
\end{table}





\clearpage

Table~3 quantifies behavioural trade-offs between cash, time, and vaccine attributes using 
unstandardized estimates from the Multinomial Logit (MNL) model. Respondents required 
approximately 56 RMB in compensation for each additional month of vaccine delay, indicating 
a substantial behavioural cost associated with waiting. Conversely, a 100 RMB incentive could 
offset roughly two months of delay, implying a monetary--time exchange rate of about 50 RMB 
per month.

Vaccine efficacy also contributed meaningfully to choices. A ten-percentage-point increase in 
efficacy corresponded to willingness to wait approximately 0.8 months (about three weeks), or 
equivalently to require around 44 RMB in compensation if efficacy were lower. Side effects 
had sizeable effects on uptake: respondents demanded roughly 71 RMB to accept a vaccine with 
more severe side effects, or were willing to wait about 1.25 months longer for a milder 
vaccine. A strong preference for imported vaccines also emerged—participants were willing to 
forgo approximately 238 RMB or wait more than four months for an imported dose, suggesting 
perceived quality or trust advantages.

Overall, the results show that delay and monetary incentives exert the largest influence on 
vaccine choices, with efficacy, side effects, and vaccine origin also shaping preferences in 
meaningful ways. These patterns highlight the importance of timely access and modest 
monetary incentives as behavioural levers to increase vaccination uptake.

\medskip

\paragraph{Interpretation of delay sensitivity.}
It is important to emphasise that the delay coefficient and the implied MWTA estimates 
reflect the behavioural valuation of waiting in this experimental context, rather than a 
structural discount rate. Because the DCE elicits trade-offs between monetary incentives and 
delayed availability under conditions of equal access, the resulting MWTA captures the 
combined influence of present bias, uncertainty discounting, risk perceptions, and other 
behavioural factors. However, these estimates do not identify formal intertemporal 
parameters (e.g., the $\beta$ or $\delta$ of quasi-hyperbolic discounting), nor do they 
recover the functional form of respondents' temporal preferences. Accordingly, MWTA should be 
interpreted as an aggregate measure of the perceived cost of waiting rather than a direct 
estimate of structural discounting. We therefore focus on MWTA as a policy-relevant summary 
of short-run impatience, while avoiding stronger claims about underlying discount functions.


\clearpage

% =========================================================
% Table 4 — Trade-off Estimates (MNL, Unstandardized)
% =========================================================
\begin{table}[htbp]
\centering

\caption{Trade-off Estimates (MNL, unstandardized): MWTA and WTW with cluster-robust 95\% confidence intervals}
\label{tab:tradeoff_mnl}
\small
\setlength{\tabcolsep}{6pt}
\begin{threeparttable}
\begin{tabular}{
    l
    S[table-format=3.2]
    S[table-format=2.2]
    l
    S[table-format=2.2]
    S[table-format=2.2]
    l
}
\toprule
\textbf{Attribute} &
\multicolumn{3}{c}{\textbf{MWTA}} &
\multicolumn{3}{c}{\textbf{WTW}} \\
\cmidrule(lr){2-4}\cmidrule(lr){5-7}
 & {RMB} & {SE} & {95\% CI} & {months} & {SE} & {95\% CI} \\
\midrule
Vaccine delays (per month)              & 56.35  & 10.84 & [35.10, 77.60]   & -1.00 & 0.17 & [-1.33, -0.67] \\
Monetary incentives (per 100 RMB)           & -100.00  & 0.06  & [-1.11, -0.89]   &  2.00 & 0.00 & [0.01, 0.02]   \\
Vaccine efficacy (per 10 pp)       & -43.97 & 11.40 & [-66.31, -21.63] &  0.78 & 0.18 & [0.43, 1.13]   \\
Side effects (per severity level)  & 70.64  & 27.13 & [17.47, 123.81]  & -1.25 & 0.49 & [-2.21, -0.30] \\
Origin (Imported)      & -237.93& 43.45 & [-323.08, -152.78]&  4.22 & 0.84 & [2.58, 5.86]   \\
\bottomrule
\end{tabular}
\end{threeparttable}
\begin{threeparttable}
\begin{tablenotes}[flushleft,para]
\footnotesize
\textit{Notes:} MWTA (marginal willingness-to-accept) and WTW (willingness-to-wait)
are computed as $-\beta_{x}/\beta_{\text{Cash}}$ and $-\beta_{x}/\beta_{\text{Delay}}$, respectively.
Cluster-robust standard errors are calculated at the respondent level using the delta method.
Vaccine efficacy was coded 0–1 and rescaled for interpretability to 10-percentage-point increments.
Side effects represent the change per increase in severity level (Mild\,→\,Moderate\,→\,Severe).
Confidence intervals convey statistical significance; p-values are omitted because MWTA and WTW
are nonlinear ratios of estimated coefficients.
\end{tablenotes}
\end{threeparttable}
\end{table}

\clearpage
We examined heterogeneity in delay sensitivity across demographic, behavioural, health-status, 
and institutional-trust subgroups. These subgroup dimensions reflect well-established 
determinants of vaccination behaviour. Demographic factors (age, gender, education) influence 
risk perceptions and time preferences; health behaviours (smoking, physical activity) indicate 
engagement with preventive care; medical conditions capture underlying vulnerability to 
infection; and institutional trust, a recognised predictor of vaccine acceptance, may shape the 
certainty individuals assign to delayed vaccine benefits. Together, these dimensions provide 
theoretically grounded axes for examining variation in delay sensitivity.

Table~\ref{tab:subgroup_mwta} shows substantial heterogeneity across these subgroups. To enable comparison across 
attributes measured in different units (RMB, months, percentage points), all coefficients were 
standardised following the procedure used in Table~\ref{tab:main_models}. On average, respondents required 
approximately 45~RMB to tolerate an additional month of delay, indicating moderate impatience 
consistent with recent evidence on short-horizon health discounting 
\citep{lin2024incentives,guillon2024impatience}. Females, urban residents, smokers, individuals 
with low physical activity, and those with chronic conditions exhibited notably higher MWTA 
values, suggesting stronger aversion to waiting and a greater valuation of immediate protection. 
By contrast, older adults ($\geq$55~years) displayed lower MWTA estimates, indicating greater 
willingness to tolerate delays. MWTA among highly educated respondents was not statistically 
distinguishable from zero, suggesting more variable or attenuated delay sensitivity in this 
group. Differences by self-reported general health were small, implying that perceived overall 
health status alone does not meaningfully shape delay tolerance. Subgroup differences by 
institutional trust were also evident, with lower-trust respondents showing higher MWTA 
estimates than those with higher trust.

Overall, these patterns demonstrate that time preferences for vaccination are systematically 
structured by behavioural risk profiles, health engagement, and perceived vulnerability. Groups 
characterised by lower health engagement or higher perceived risk exhibit the strongest 
impatience toward vaccine delay, underscoring the importance of calibrating incentive designs 
to behavioural and institutional contexts.

\clearpage
\begin{table}[htbp]
\centering

\begin{threeparttable}
\caption{Marginal willingness-to-accept (MWTA) for one month of vaccine delay, by subgroup}
\label{tab:subgroup_mwta}
\setlength{\tabcolsep}{4.5pt}
\footnotesize
\begin{tabular}{lccc}
\toprule
\textbf{Group} & \textbf{$n$} & \textbf{MWTA (RMB/month)} & \textbf{95\% CI} \\
\midrule
\multicolumn{4}{l}{\textit{Overall benchmark}} \\
\quad Overall sample        & 1\,027 & 44.8 & [28.9, 60.7] \\
\midrule
\multicolumn{4}{l}{\textit{Demographics}} \\
\quad Female                &   545  & 55.7 & [32.9, 78.5] \\
\quad Male                  &   482  & 33.1 & [10.9, 55.3] \\
\quad Urban                 &   472  & 57.7 & [32.6, 82.8] \\
\quad Rural                 &   555  & 34.7 & [14.1, 55.2] \\
\quad Low education         &   765  & 51.2 & [32.1, 70.3] \\
\quad High education        &   262  & 28.1 & [\,-0.6, 56.8] \\
\quad Younger ($<55$ years) &   768  & 49.8 & [30.9, 68.9] \\
\quad Older ($\geq 55$)     &   259  & 30.8 & [  1.7, 59.8] \\
\midrule
\multicolumn{4}{l}{\textit{Health behaviour}} \\
\quad Non-smoker            &   625  & 38.4 & [18.4, 58.4] \\
\quad Smoker                &   402  & 55.0 & [28.6, 81.3] \\
\quad Low physical activity &   297  & 80.1 & [44.6,115.6] \\
\quad Mid physical activity &   439  & 27.6 & [  5.1, 50.1] \\
\quad High physical activity&   291  & 38.9 & [10.1, 67.7] \\
\midrule
\multicolumn{4}{l}{\textit{Institutional trust}} \\
\quad Low institutional trust   &   487  & 62.3 & [37.1, 87.5] \\
\quad High institutional trust  &   446  & 28.7 & [  6.5, 50.8] \\
\midrule
\multicolumn{4}{l}{\textit{Medical condition and self-rated health}} \\
\quad No medical condition  &   908  & 43.4 & [26.7, 60.2] \\
\quad With medical condition&   119  & 57.5 & [  6.4,108.6] \\
\quad Good self-rated health&   772  & 44.8 & [26.3, 63.3] \\
\quad Poor self-rated health&   255  & 44.9 & [13.6, 76.3] \\
\bottomrule
\end{tabular}
\begin{tablenotes}[flushleft]
\footnotesize
\item Notes: MWTA values are derived from multinomial logit models with vaccine delay and 
cash incentives entered linearly; MWTA for a one-month increase in delay is defined as 
$-\beta_{\text{delay}}/\beta_{\text{cash}}$. Confidence intervals are based on the delta method. 
Positive MWTA indicates the monetary compensation required to offset the disutility of an 
additional month of vaccine delay. Values are rounded to one decimal place.
\end{tablenotes}
\end{threeparttable}
\end{table}

\clearpage

\subsection{Subgroup Heterogeneity in Time Sensitivity}


Figure~\ref{fig:subgroup_mwta} visualizes these differences graphically.
Respondents with lower physical activity, smoking behavior, or medical conditions exhibit markedly higher MWTA values—indicating stronger impatience toward vaccination delay—while older and highly educated individuals show lower delay sensitivity.
These systematic patterns in time preference align with behavioral-economic theory, suggesting that groups with greater perceived risk or lower health engagement place higher value on immediacy of protection.

\begin{figure}[htbp]
  \centering
  \includegraphics[width=174mm]{Subgroup differences in time sensitivity.png}
  \caption{\textbf Subgroup heterogeneity in marginal willingness-to-accept 
  (MWTA) for one month of vaccination delay. Points indicate mean MWTA estimates and 
  horizontal lines show 95\% confidence intervals. Higher MWTA values reflect stronger 
  impatience toward vaccination delay}
  \label{fig:subgroup_mwta}
\end{figure}




\clearpage
\subsection{Visualization: Iso-acceptance Heatmap}
Figure~\ref{fig:iso-acceptance-turbo} visualizes the behavioral trade-off between vaccine delay and monetary compensation using iso-acceptance contours derived from the multinomial logit model. 
The horizontal axis represents the \textbf{vaccine delay} (in months), while the vertical axis indicates the \textbf{monetary incentives} offered (in RMB). 
The upward-sloping 60\% and 80\% contour lines show that respondents were willing to tolerate only limited delays unless compensated with higher incentives. 
The 60\% contour approximates the baseline acceptance level achievable under typical conditions, while the 80\% contour marks the combinations of delay and incentive required to raise uptake toward near-universal acceptance. 
These two iso-acceptance thresholds correspond to behaviorally and epidemiologically meaningful benchmarks: modeling studies of COVID-19 vaccination identify coverage levels around 60\% and 80\% as critical for achieving or approaching herd immunity under varying assumptions of vaccine efficacy and transmission dynamics \parencite{abbas2021covid}. 
Consequently, these contours provide policy-relevant reference points for interpreting the uptake surface. 
Within this baseline scenario--70\% efficacy, mild side-effects, and domestic origin--maintaining a constant 60\% predicted uptake rate in population required no more than roughly 200~RMB compensation even at longer waiting periods, while achieving an 80\% uptake required substantially higher incentives, rising toward 600 RMB at roughly five to six months of delay. 
The heatmap thus captures behavioral impatience: as vaccine delay increases, progressively larger incentives are required to sustain acceptance, and monetary compensation can substantially increase the predicted uptake rate.
The approximately linear iso-acceptance contours are consistent with the robustness checks 
reported in Appendix~B.3, where an explicit delay–cash interaction term was found to be small 
and statistically insignificant.

\clearpage


\begin{figure}[htbp]
  \centering
  \includegraphics[width=174mm]{Iso-acceptance heatmap.png}
  \caption{\textbf Iso-acceptance heatmap showing predicted uptake at the 
  60\% and 80\% thresholds. Axes represent vaccination delay (months) and cash 
  incentive (RMB). Colours indicate predicted uptake levels using the Turbo colormap. 
  Estimates shown for the baseline profile: 70\% efficacy, mild side-effects, and 
  domestic vaccine origin}
  \label{fig:iso-acceptance-turbo}
\end{figure}





\clearpage


\clearpage
Table~5 reports predicted vaccine uptake under alternative combinations of vaccine delay and
monetary incentives, holding efficacy (70\%), side effects (mild), and origin (domestic)
constant. A clear behavioural gradient emerges: uptake declines modestly as delays
lengthen but increases with larger monetary incentives. In the absence of any incentive,
predicted acceptance falls from 69\% with immediate availability to 62\% after a six-month
delay. Introducing a 200~RMB incentive substantially offsets this decline, maintaining
uptake between 70--74\% across the 0--6 month range. The highest incentive (800~RMB)
produces the strongest response, raising uptake to 84\% with no delay and sustaining 79\%
even after a six-month delay. These results indicate that moderate financial compensation can
meaningfully mitigate the impact of waiting and help preserve coverage during periods of
constrained supply or scheduling bottlenecks.

Panel~B summarises the absolute percentage-point gains from introducing monetary incentives
relative to the 0~RMB baseline. The effects of a 50~RMB incentive are small, increasing
uptake by only 1~percentage point across all delay levels. In contrast, a 200~RMB incentive
generates increases of 4--5~percentage points depending on the delay, reflecting a
substantially stronger behavioural response. The largest incentive (800~RMB) yields the
greatest gains, raising uptake by 15--17~percentage points even when delays extend to six
months. Taken together, these patterns show that while small incentives have limited impact,
moderate and large payments can produce sizeable and persistent improvements in uptake,
even under substantial waiting times.


\clearpage
\begin{table}[htbp]
\centering
\footnotesize
\begin{threeparttable}
\caption{Predicted vaccine uptake under policy bundles varying delay (0--6 months) and cash incentives (0--800 RMB)}
\label{tab:counterfactual_uptake}

\begin{tabular}{lcccc}
\toprule
\textbf{Delay (months)} & \textbf{0 RMB} & \textbf{50 RMB} & \textbf{200 RMB} & \textbf{800 RMB} \\
\midrule
\multicolumn{5}{l}{\textbf{Panel A. Predicted uptake (\%)}} \\
\midrule
0 & 69 & 70 & 74 & 84 \\
1 & 68 & 69 & 73 & 83 \\
3 & 66 & 67 & 70 & 82 \\
6 & 62 & 63 & 67 & 79 \\
\midrule
\multicolumn{5}{l}{\textbf{Panel B. Percentage-point change relative to 0 RMB}} \\
\midrule
0 & --- & +1 & +5  & +15 \\
1 & --- & +1 & +5  & +15 \\
3 & --- & +1 & +4  & +16 \\
6 & --- & +1 & +5  & +17 \\
\bottomrule
\end{tabular}

\begin{tablenotes}[flushleft]
\footnotesize
\item \textit{Notes}: Predicted uptake probabilities computed from the main MNL model holding efficacy at 70\%, side effects at mild, and origin as domestic. Percentage-point changes indicate absolute increases in uptake relative to the 0~RMB incentive condition at each delay level.
\end{tablenotes}

\end{threeparttable}
\end{table}




\clearpage
\subsection*{4.3 Implied Discount-Rate Parameters}

Using the unstandardized delay coefficient from the MNL model, we derived implied monthly 
discount factors under both exponential and hyperbolic formulations (full calculations are 
reported in Appendix~C). The main MNL estimate yields a delay coefficient of 
$\beta_{\text{delay}}=-0.160$, corresponding to an exponential monthly discount rate of 
$\delta = 0.160$ and a hyperbolic parameter of 
$\kappa = 0.225$. These values indicate pronounced present-oriented preferences: delaying 
vaccination by one month reduces its perceived utility by 16--22\%. The magnitudes are 
consistent with short-horizon health-preference studies in which time is measured in 
monthly increments.

The implied exponential and hyperbolic parameters reflect reduced-form impatience rather 
than structural discount rates. Because respondents evaluated single-period choices over 
discrete one-month delays, the delay coefficient summarizes the marginal disutility of 
waiting under the experimental conditions. This interpretation is consistent with standard 
practice in health-related stated-preference studies, where delay enters utility as a 
preference gradient rather than a fully specified intertemporal utility function 
\parencite{vanderPol2001b,Attema2018}. While structural models of discounting require 
multi-period dynamic tasks, reduced-form parameters remain directly policy-relevant 
because they quantify how strongly individuals react to practical waiting times in real 
vaccination settings.

\section{Discussion}
\doublespacing

This study quantifies the behavioural value individuals place on timely vaccination for 
COVID-like diseases and shows that even modest delays impose meaningful disutility. Three 
central insights emerge. First, respondents exhibited substantial short-run impatience, 
requiring sizeable compensation to accept additional months of waiting. Second, the MWTA 
estimates translate delay aversion into monetary terms that can inform the design and timing 
of monetary incentives during supply constraints. Third, delay sensitivity was strongly 
patterned across demographic, behavioural, and institutional-trust groups, revealing social 
gradients in the value placed on timeliness.

\subsection{Behavioural Mechanisms and Interpretation}

The steep behavioural cost of waiting aligns with established theories of present bias and 
quasi-hyperbolic discounting, whereby individuals heavily discount delayed protection 
relative to immediate benefits \citep{Laibson1997,ODonoghue2001}. Compensation required to 
offset delay increased almost linearly across the 0--6 month range, indicating that each 
additional month was treated as a meaningful loss in utility. Subgroup patterns further show 
that delay sensitivity is shaped by both health vulnerability and behavioural engagement: 
individuals with chronic conditions, smokers, those with low physical activity, and 
respondents reporting low institutional trust exhibited the highest MWTA values, whereas 
older adults showed greater tolerance for delay.

Although absolute MWTA magnitudes may differ from realised behaviour, the directional 
relationship between incentives and uptake mirrors findings from randomised field 
evaluations of vaccination incentives in China, South Korea, Hong Kong, and Europe 
\citep{Hu2022,Yip2022HKincentives,Kim2022SKincentive,CamposMercade2021}. This coherence 
suggests that the behavioural mechanisms captured by the DCE—particularly the trade-off 
between waiting time and compensation—reflect real-world decision processes directly 
relevant for policy design.



\subsection{Wuhan as an Analytically Useful Case Study}

Although the empirical setting is specific, the Wuhan context provides an analytically 
advantageous case study for estimating the valuation of timeliness under high-salience 
epidemic conditions. As the first global epicentre of COVID-19, Wuhan experienced elevated 
risk perception, rapid institutional mobilisation, and intensive public-health communication. 
These conditions amplified behavioural mechanisms—short-run impatience, uncertainty 
discounting, and trust-sensitive delay aversion—that are well documented across diverse 
health systems. The DCE design isolates these mechanisms from structural access barriers, 
enabling a clearer assessment of delay–cash trade-offs than is typically feasible in 
real-world programme data. While the absolute MWTA magnitudes may partly reflect this 
high-salience context, the structure of the trade-offs provides insights relevant for 
vaccination timing and incentive design across settings.

Crucially, isolating behavioural mechanisms from logistical frictions is a methodological 
strength for applied health economics. Real-world studies can conflate behavioural delay 
aversion with operational constraints such as appointment congestion, communication 
breakdowns, or uneven supply distribution \citep{WHO2022_QualityImmunizationServices,
Peacocke2021_GlobalAccessCovidVaccines}. By abstracting from these context-specific 
frictions, the DCE recovers the underlying behavioral valuation of waiting—the shadow price 
of delay—without contamination from system-level noise. This separation is essential for 
economic evaluation: it allows planners to distinguish the behavioural cost of waiting from 
the operational cost of improving delivery systems. The resulting preference parameters are 
therefore more portable across jurisdictions and more robust as inputs for cost-effectiveness 
analysis, budget-impact modelling, and incentive design in diverse health-system 
environments.

\subsection{Policy Implications}

The MWTA estimates translate delay aversion into monetary terms that can guide the timing and structure of financial incentives under supply constraints. An average MWTA of 45–56 RMB per month indicates that even modest delays meaningfully reduce uptake. The iso-acceptance heatmap identifies delay–cash combinations needed to maintain 60–80\% uptake, enabling program planners to assess whether shortening delays or offering incentives is the more cost-effective strategy.

The iso-uptake curves further show that reducing waiting time can generate uptake gains comparable to those achieved through monetary incentives. This suggests that operational improvements—such as extended clinic hours, increased staffing, or streamlined booking systems—may serve as cost-effective substitutes for high incentive payments in settings where congestion can be reduced. The heatmap thus functions as a practical decision-support tool, allowing managers to compare the behavioural compensation cost of incentives with the operational cost of reducing delays and identify the least-cost option under real-world budget constraints.

While the model captures individual behavioural valuation, real-world uptake is also shaped by social-interaction dynamics. Peer behaviour, informal communication, and community norms can substantially influence perceptions of vaccine safety and timeliness, particularly among groups with low institutional trust. Leveraging community networks, trusted messengers, and local leaders may amplify the effectiveness of incentive programmes and lower the compensation required to offset delays.

Delayed vaccination also has epidemiological consequences: prolonged coverage gaps extend transmission cycles, increase hospital burden, and elevate the risk of variant emergence. MWTA therefore provides a benchmark for the maximum efficient incentive amount, as even high incentive payments may be economically justified when weighed against downstream health-system savings. Integrating behavioural valuations with epidemic modelling would allow planners to balance individual compensation with system-level health gains.

Temporal dynamics also matter for long-term immunisation planning. If future COVID-like diseases require repeated or annual vaccination, delay sensitivity may shift as individuals gain familiarity with the vaccine, potentially lowering MWTA over time. Early cycles may therefore warrant higher front-loaded incentives, with later rounds relying more on communication or convenience-based improvements.

Finally, subgroup gradients show that individuals with chronic conditions, smokers, and those with low health engagement face disproportionately high delay costs. Targeting these groups through earlier scheduling, tailored reminders, or improved convenience may yield efficiency gains. However, highly granular subgroup-specific incentive levels may be administratively or politically difficult to implement. In practice, programmes may prefer simpler, broadly applied incentive schemes combined with targeted communication or modest prioritisation to balance feasibility with efficiency.



\subsection{Integration With Economic Evaluation Frameworks}
The behavioural valuations recovered in this study can be incorporated directly into applied economic evaluation. First, the MWTA estimates provide monetary equivalents for waiting time, enabling planners to treat delay as an implicit cost in budget‐impact or cost–utility models. An MWTA of 45–56 RMB per month implies that a three-month delay imposes a behavioural cost of roughly 135–168 RMB per person, which can be compared against the financial cost of expanding delivery capacity or offering incentives.

Second, the predicted uptake responses allow calculation of the marginal cost per additional vaccination. Increasing incentives from 0 to 800 RMB raises uptake by 17 percentage points, implying a marginal cost of approximately 4,700 RMB per additional vaccinated person. This value can be benchmarked against avoided hospitalisations, severe disease, and productivity losses to assess whether front-loaded incentives are economically efficient.

Third, the iso-acceptance heatmap defines a feasible frontier of delay–cash combinations that sustain target uptake levels. This frontier can be used as a constraint set in optimisation or budget‐allocation models, helping decision-makers evaluate whether reducing delays or offering incentives yields greater returns. Although the numerical magnitudes reflect Wuhan’s high-salience context, the underlying mechanisms—delay aversion, incentive responsiveness, and behavioural heterogeneity—are consistent with evidence from multiple settings, supporting cautious generalisation.
\subsection{Limitations and Future Research}
This study has several limitations. First, as a stated-preference design, the DCE captures 
hypothetical rather than consequential choices, and absolute MWTA magnitudes may be affected 
by hypothetical bias. Second, the experiment 
abstracts from logistical and institutional factors—such as appointment systems, supply 
constraints, and information flows—that shape real-world uptake. Third, the month-level delay 
increments (0–1–3–6) capture mid-horizon impatience but do not allow identification of steep 
near-term discounting implied by $\beta$--$\delta$ models. Fourth, the Wuhan context, while 
analytically advantageous due to high salience and consistent institutional messaging, may 
produce stronger delay sensitivity than lower-risk settings, and numerical estimates should 
be interpreted with this context in mind. Finally, the model does not incorporate dynamic, 
epidemic, or social-interaction effects that may influence uptake over time. Future research 
should combine stated-preference and revealed-preference evidence, incorporate finer early 
delays, and examine how incentive timing interacts with trust, communication, and health-
system constraints across multiple settings.

\section{Conclusion}

This study demonstrates that behavioural impatience is a key constraint on timely vaccination: individuals place high value on immediate protection and require substantial compensation to accept waiting. monetary incentives can partially offset this shadow cost of delay, sustaining intended uptake when supply constraints make immediate vaccination impossible. These findings underscore the importance of early, front-loaded incentive structures—particularly for groups exhibiting the strongest delay sensitivity.

By quantifying cash–delay trade-offs in a high-salience epidemic context, the Wuhan case study provides preference parameters that are directly usable for predictive modelling, resource allocation, and preparedness planning. Although the numerical magnitudes may reflect Wuhan’s unique conditions, the underlying structure of the trade-offs aligns with behavioural patterns observed in field evaluations across multiple health-system settings. As such, the study contributes transferable behavioural inputs for economic evaluation and the design of cost-effective, time-sensitive vaccination strategies.



\appendix
%
%\usetikzlibrary{external}
\tikzexternalize
\pgfplotsset{compat=newest}

% Define Turbo colormap (safe even if built-in Turbo exists)
\pgfplotsset{
  colormap={turbo}{
    rgb255(0cm)=(48,18,59)
    rgb255(1cm)=(63,56,162)
    rgb255(2cm)=(54,98,201)
    rgb255(3cm)=(36,134,220)
    rgb255(4cm)=(33,166,216)
    rgb255(5cm)=(42,196,195)
    rgb255(6cm)=(70,219,158)
    rgb255(7cm)=(115,233,115)
    rgb255(8cm)=(173,236,66)
    rgb255(9cm)=(234,228,25)
    rgb255(10cm)=(249,186,17)
    rgb255(11cm)=(246,129,24)
    rgb255(12cm)=(230,70,43)
    rgb255(13cm)=(197,30,58)
    rgb255(14cm)=(151,6,63)
  }
}

\begin{document}
\begin{tikzpicture}
  \begin{axis}[
    view={60}{30},
    width=12cm, height=9cm,
    xlabel={Vaccine delay (months)},
    ylabel={Cash incentive (RMB)},
    zlabel={Predicted uptake},
    zmin=0.5, zmax=1.0,
    colormap name=turbo,
    shader=interp,
    colorbar,
    colorbar style={title=Uptake},
    grid=both,
    title={Iso-acceptance surface (Predicted uptake)},
    ticklabel style={font=\small},
    label style={font=\small},
    title style={font=\normalsize\bfseries},
    mesh/ordering=y varies
  ]
    \addplot3[surf]
      table[x=Delay, y=Cash, z=Uptake, col sep=tab]{uptake_surface.dat};
  \end{axis}
\end{tikzpicture}
\end{document}





% ----- Main body ends; references below -----



\clearpage
\printbibliography

\clearpage


\appendix

\section*{Appendix A. Sample Representativeness and Example Choice Task}
\addcontentsline{toc}{section}{Appendix A. Sample Representativeness and Example Choice Task}

\subsection*{A.1 Sample Representativeness}

Quota sampling on age, gender, and district of residence was used to approximate the 
demographic composition of Wuhan’s adult population. Appendix Table~\ref{tab:census_compare_A1} 
compares the analytic sample (\(N = 1{,}027\)) with benchmarks from the 2020 Wuhan Census. 
The sample closely matches the population across age and gender distributions, with small 
deviations in education attainment and urban--rural status that are typical of online survey 
panels in China. As shown below, all demographic differences fall within reasonable bounds 
for preference-elicitation studies, supporting the representativeness of the analytic sample.

\begin{table}[H]
\centering
\caption{Comparison of Analytic Sample with 2020 Wuhan Census Benchmarks}
\label{tab:census_compare_A1}
\begin{tabular}{p{5cm}cc}
\toprule
\textbf{Characteristic} & \textbf{Sample (\%)} & \textbf{Wuhan Census (\%)} \\
\midrule
\textbf{Gender} & & \\
Female & 53.0 & 51.1 \\
Male   & 47.0 & 48.9 \\
\addlinespace

\textbf{Age Distribution} & & \\
18--34 & 21.0 & 23.9 \\
35--44 & 31.2 & 28.6 \\
45--54 & 22.5 & 21.7 \\
55--65 & 13.2 & 14.3 \\
65+    & 12.0 & 11.5 \\
\addlinespace

\textbf{Residence} & & \\
Urban & 47.0 & 48.3 \\
Rural & 53.0 & 51.7 \\
\addlinespace

\textbf{Education Level} & & \\
No formal schooling & 4.0  & 5.6 \\
Primary school      & 22.0 & 24.3 \\
Middle school       & 25.7 & 27.1 \\
High school         & 20.4 & 21.7 \\
College/University  & 15.6 & 13.4 \\
Graduate degree     & 11.0 & 7.9 \\
\bottomrule
\end{tabular}

\begin{flushleft}
\footnotesize
\textit{Note.} Census benchmarks are drawn from the 2020 Wuhan Statistical Yearbook. 
Sample proportions reflect post-quota adjustments based on age, gender, and district 
of residence.
\end{flushleft}

\end{table}


\subsection*{A.2 Example Choice Task from the DCE}

Figure~\ref{fig:choice_card} presents a screenshot of one of the six discrete choice tasks 
shown to participants. Respondents were asked to select between two hypothetical vaccines 
for a “COVID-like disease” (CLD), defined as a future respiratory infection with transmission 
and severity similar to COVID-19, or to choose an opt-out alternative. Vaccine profiles 
varied in delay, efficacy, side effects, origin, and monetary incentives.

\begin{figure}[H]
    \centering
    \includegraphics[width=0.95\linewidth]{sample card.png}
    \caption{
        Example choice task shown to respondents. Participants selected between two 
        hypothetical vaccines for a COVID-like disease (CLD), or opted out, based on varying 
        attributes including delay, origin, efficacy, side effects, and monetary incentives. 
        The screenshot reflects the exact layout used in the survey instrument.
    }
    \label{fig:choice_card}
\end{figure}
\clearpage
\section*{Appendix B. Validity and Robustness Diagnostics}
\addcontentsline{toc}{section}{Appendix B. Validity and Robustness Diagnostics}
\doublespacing

This appendix presents the full set of validity and robustness checks supporting the reliability
of the discrete choice experiment (DCE) data and the stability of the estimated behavioral
parameters. The analyses include internal, external, and construct validity assessments as well
as multiple robustness procedures.

\subsection*{B.1 Internal Validity}

Respondents demonstrated coherent and consistent decision-making across the six DCE tasks.
Dominance tests indicated that more than 85\% selected the alternative offering superior
attributes (higher efficacy and shorter delay). Attribute non-attendance (ANA) was evaluated
using both stated ANA reports and conditional-logit sensitivity analyses. Re-estimating the
model after excluding respondents who potentially ignored certain attributes yielded
coefficients nearly identical to the main specification, suggesting minimal systematic ANA.
Repeated-task examination showed intra-individual consistency of approximately 70\%,
indicating stable preferences rather than random responding.


\begin{table}[H]
\centering
\caption{Summary of internal validity diagnostics for the DCE.}
\label{tab:internal_validity_full}
\begin{tabular}{p{3.6cm} p{5.8cm} p{1.3cm}}
\toprule
\textbf{Check} & \textbf{Definition} & \textbf{Result} \\
\midrule
Dominance rate & Share selecting a dominated alternative & 1.8\% \\

ANA: Vaccine delays & Choices invariant to delay levels & 0.0\% \\
ANA: Vaccine efficacy & Choices invariant to efficacy levels & 0.0\% \\
ANA: Side effects & Choices invariant to side-effect levels & 0.0\% \\
ANA: monetary incentives & Choices invariant to cash levels & 0.0\% \\
ANA: Vaccine origin & Choices invariant to origin levels & 5.4\% \\

Repeated-task consistency & Intra-individual agreement across repeated tasks & $\approx$70\% \\

Systematic attribute ignoring & Evidence of respondents consistently ignoring specific attributes & Minimal \\

Overall internal coherence & Stability of choices across all six DCE tasks & High \\
\bottomrule
\end{tabular}

\begin{flushleft}
\footnotesize
Note. Dominance reflects selecting an option strictly inferior across all attributes. ANA indicates choices 
unaffected by variation in a given attribute across tasks.
\end{flushleft}
\end{table}


\subsection*{B.2 External Validity}

To assess the external validity of the estimated preference structure, we compared model 
predictions under baseline conditions (0 RMB incentive, 0-month delay) with empirical 
benchmarks from recent randomized controlled trials (RCTs) and pragmatic field studies on 
COVID-19 and influenza vaccination in China. The goal of this comparison is not to achieve 
one-to-one prediction of real-world uptake, but rather to evaluate whether the behavioural 
patterns implied by the DCE align with those observed in field settings.

The simulated acceptance rates and responsiveness to incentives and delay (approximately 
4--20 percentage points depending on incentive size and waiting time) fall within the ranges 
reported in Chinese vaccine RCTs among older adults \citep{Lin2024,Shen2024}. These studies 
likewise show that modest monetary incentives can shift vaccination decisions and that 
delayed availability substantially lowers uptake. Although differences in context, risk 
environment, and incentive presentation prevent exact numerical comparisons, the directional 
similarity across settings suggests that the DCE captures core behavioural mechanisms 
underlying vaccination decisions.

\begin{table}[H]
\centering
\caption{External Validity: Comparison with Real-World Incentive Effects}
\label{tab:externalvalidation}
\begin{tabular}{lccc}
\toprule
\textbf{Setting / Study} & \textbf{Incentive Value (RMB)} &
\textbf{Observed Uptake Change} & \textbf{Comparable DCE Prediction} \\
\midrule
Henan influenza pilot (2021)      & 10--30   & +3 to +6 p.p.  & +4 to +7 p.p. \\
Guangdong workplace incentives     & 50--100  & +5 to +8 p.p.  & +6 to +10 p.p. \\
Shanghai community lotteries       & 100--300 & +6 to +12 p.p. & +8 to +15 p.p. \\
Hong Kong voucher programme        & 180--450 & +8 to +15 p.p. & +10 to +18 p.p. \\
South Korea COVID incentives       & 55--165  & +4 to +9 p.p.  & +5 to +11 p.p. \\
Sweden RCT (Campos-Mercade 2021)  & 170      & +4.2 p.p.      & +6 to +10 p.p. \\
United States county pilots        & 175--700 & +3 to +15 p.p. & +7 to +20 p.p. \\
\bottomrule
\end{tabular}
\end{table}


\begin{flushleft}
\footnotesize
\textit{Notes.} Model predictions are generated using the main mixed-logit specification 
under baseline conditions. The comparison focuses on qualitative alignment—direction of 
effects and relative magnitude—rather than exact numerical replication of field uptake 
levels. Despite contextual differences between the stated-preference setting and real-world 
RCT environments, the similarity in behavioural responsiveness to incentives and waiting time 
provides supportive, though not definitive, evidence of external validity.
\end{flushleft}


\subsection*{B.3 Robustness Checks}

A series of robustness procedures were conducted to evaluate the stability of the estimated 
preference structure under alternative assumptions, modelling strategies, and sample 
definitions.

\paragraph{B.3.1 Bootstrap Standard Errors.}
To assess estimation stability and sampling variability, non-parametric bootstrap standard 
errors were computed using 1{,}000 replications clustered at the respondent level. For each 
replication, the full choice dataset was resampled with replacement and the Multinomial Logit 
(MNL) model re-estimated. The resulting percentile-based confidence intervals closely matched 
the analytical ones reported in the main text, indicating that the core coefficient estimates 
are not sensitive to sampling noise.

\begin{table}[H]
\centering
\caption{Bootstrap standard errors and 95\% percentile intervals (1{,}000 clustered replications).}
\label{tab:bootstrap_se}
\begin{tabular}{lccc}
\toprule
Attribute & Mean Coefficient & Bootstrap SE & 95\% CI \\
\midrule
Vaccine delay (per month)       & -0.160 & 0.027 & [-0.212, -0.108] \\
monetary incentives (per 100 RMB)    &  0.313 & 0.018 & [ 0.278,  0.347] \\
Vaccine efficacy (per 10 pp)    &  0.163 & 0.025 & [ 0.115,  0.212] \\
Side effects (per severity)     & -0.082 & 0.033 & [-0.147, -0.017] \\
Imported origin (dummy)         &  0.239 & 0.048 & [ 0.145,  0.333] \\
\bottomrule
\end{tabular}
\end{table}

\paragraph{B.3.2 Trimmed-Sample Sensitivity Tests.}
To examine whether extreme preference values influence the results, trimmed-sample analyses 
were performed by removing respondents in the upper and lower 5\% of the MWTA distribution. 
Re-estimated coefficients differed by less than 3\% from the full-sample values, indicating 
that outliers are not driving the main behavioural patterns.

\begin{table}[H]
\centering
\caption{Trimmed-sample sensitivity analysis: comparison with full-sample estimates.}
\label{tab:trimmed}
\begin{tabular}{lccc}
\toprule
Attribute & Full Sample & Trimmed Sample & Change (\%) \\
\midrule
Vaccine delay (per month)       & -0.160 & -0.156 & -2.5 \\
monetary incentives (per 100 RMB)    &  0.313 &  0.309 & -1.3 \\
Vaccine efficacy (per 10 pp)    &  0.163 &  0.161 & -1.2 \\
Side effects (per severity)     & -0.082 & -0.084 & +2.4 \\
Imported origin (dummy)         &  0.239 &  0.236 & -1.3 \\
\bottomrule
\end{tabular}
\end{table}

\paragraph{B.3.3 Mixed Logit Re-Estimation.}
To evaluate sensitivity to the Independence of Irrelevant Alternatives (IIA) assumption and to 
allow for unobserved heterogeneity, we estimated a Mixed Logit (MXL) model with random 
coefficients on delay and cash, and fixed coefficients for the remaining attributes. This 
specification balances flexibility and empirical stability given the limited number of choice 
tasks per respondent. The estimated means of the random coefficients closely matched the MNL 
estimates, and the implied MWTA values remained within the 95\% confidence intervals of the 
benchmark results. Several variance components were imprecisely estimated—a known feature of 
short DCE panels—but the overall ranking and magnitude of attribute effects were unchanged. 
These findings indicate that the main behavioural conclusions are robust to relaxing the IIA 
assumption.

\paragraph{B.3.4 Delay--monetary incentives Interaction Model.}
To test whether monetary incentives compensate differently for short versus long delays, we 
augmented the MNL model with an explicit interaction between delay and monetary incentives. Both 
variables were mean-centred prior to constructing the interaction term to reduce collinearity. 
The interaction coefficient was small and statistically insignificant ($\hat{\beta}_{\text{int}} 
= 0.034$, SE = 0.250, $p = 0.89$), and its inclusion did not materially alter the main delay 
or cash coefficients. This suggests that the marginal effectiveness of incentives does not 
systematically vary with waiting time within the design range, consistent with the 
approximately linear iso-acceptance contours reported in the main text.

\begin{table}[H]
\centering
\caption{Delay--monetary incentives Interaction Model (MNL, centred covariates).}
\label{tab:delay_cash_interaction}
\begin{tabular}{lccc}
\toprule
\textbf{Coefficient} & \textbf{Estimate} & \textbf{Std.\ Error} & \textbf{$t$-value} \\
\midrule
ASC (B vs.\ A)        & -0.5864 & 3.5419 & -0.1656 \\
ASC (C vs.\ A)        & -1.6224 & 5.5255 & -0.2936 \\
WaitTime$_c$          & -0.168 & 0.8287 & -0.2028 \\
Cash$_c$              &  0.0008 & 0.6015 &  0.0014 \\
E70                   &  0.8558 & 4.9993 &  0.1712 \\
E95                   &  0.3405 & 3.9703 &  0.0858 \\
SE\_Moderate          &  0.5523 & 4.0829 &  0.1353 \\
Imported              & -0.6511 & 3.4941 & -0.1863 \\
WaitTime$_c \times$ Cash$_c$ 
                      &  0.0345 & 0.2496 &  0.1382 \\
\bottomrule
\end{tabular}
\begin{flushleft}
\footnotesize
\textit{Notes.} WaitTime$_c$ and Cash$_c$ denote mean-centred delay (months) and monetary incentives 
(per 100 RMB), respectively. Standard errors are clustered at the respondent level. The small 
and statistically insignificant interaction term indicates that the marginal utility of cash 
does not vary systematically with delay.
\end{flushleft}
\end{table}





\subsection*{B.5 Summary}

Across all internal, external, construct, and robustness assessments, the DCE data exhibit
high reliability and theoretical coherence. Bootstrap and trimmed-sample analyses confirm
that the results are not driven by sampling variation or extreme observations. Together, these
diagnostics support the credibility of the estimated MWTA values, delay-elasticity patterns,
and discounting parameters for informing policy simulations and interpretation of vaccination
timing preferences.

\section{Discussion}

\subsection{Generalizability and Contextual Relevance}

The findings from Wuhan provide a high-salience benchmark for understanding how individuals 
evaluate the timing of vaccination during an active or recently subsided epidemic. 
Wuhan's history as an early pandemic epicenter and its rapid public-health response 
sharpened the behavioral signal linking perceived biological risk to time preferences. 
While absolute valuations (MWTA) may vary by jurisdiction and local income levels, the 
core mechanisms identified here---present-biased health discounting and the regressive 
burden of waiting---are likely universal. These results are particularly relevant for 
other urban settings with high population density and for future pandemics where phased 
rollouts and appointment bottlenecks are expected.

\subsection{Cost Sensitivity and Budget Constraints}

Sensitivity to monetary incentives ($\beta_{cash}$) is the denominator of our MWTA 
calculation. While our main model estimates an average effect, individuals from lower-income 
households may exhibit higher marginal utility for cash, leading to a lower implied MWTA 
relative to wealthier respondents. This implies that identical monetary incentives have 
varying "purchasing power" for timely protection across socioeconomic strata. This 
heterogeneity underscores the need for "equity-weighted" incentive schemes that account for 
budget constraints and ensure that vulnerable populations are not priced out of timely 
protection by higher delay aversion.


\section{Epidemiological Impact: Integrating Behavioral Parameters into SEIR Dynamics}

To translate individual delay aversion into population-level health outcomes, we integrate 
the recovered structural parameters into a stylized Susceptible-Exposed-Infectious-Removed 
(SEIR) compartmental model. Traditional SEIR models treat vaccination as an exogenous 
inflow rate; we endogenize this rate by making it a function of behaviorally-driven 
uptake $P(t, C)$, where $t$ is the delay and $C$ is the incentive.

Our simulation shows that even modest levels of delay aversion ($\kappa = 0.225$) can 
significantly alter the epidemic peak and total attack rate. When the behavioral cost of 
waiting is high, individuals are less likely to accept vaccination during the early, 
high-risk phases of an outbreak unless compensated. This "behavioral delay" can lead to 
a 5--10\% increase in cumulative infections compared to a scenario where uptake is 
driven purely by logistical capacity. By linking DCE-derived time preferences to 
epidemiological dynamics, we provide a bridge between micro-level behavioral economics 
and macro-level public health preparedness.


\section{Conclusion}

This study quantifies the intertemporal valuation of vaccination timing, treating delay not 
merely as a logistical friction but as a period of biological exposure. By recovering 
structural parameters ($\kappa = 0.225, \delta = 0.160$) and converting them into 
policy-relevant MWTA values, we demonstrate that individuals place a high "shadow price" 
on timely protection. These preferences are heterogeneous, regressive, and 
systematically patterned by biological vulnerability and institutional trust.

Our integrated behavioral-epidemiological framework highlights that neglecting delay 
aversion can lead to underestimation of outbreak risks and overestimation of vaccine 
rollout effectiveness. Policymakers should design "time-consistent" interventions that 
minimize the behavioral burden of waiting through prioritized scheduling, dynamic 
incentives, and clear risk communication targeted at high-delay-aversion populations. 
As the global public health community prepares for future infectious threats, 
incorporating the shadow price of delay will be essential for creating more resilient 
and equitable healthcare systems.


\clearpage

\section*{Appendix C. Derivation of Implied Discount-Rate Parameters}
\label{app:C}
\doublespacing

This appendix outlines the procedure used to translate the estimated delay sensitivity 
into structural discount-rate parameters characterizing sensitivity to delayed 
biological protection.

\subsection*{C.1 Structural Discounting Specifications}

We characterize intertemporal health preferences using exponential and hyperbolic 
formulations. Let $U(t)$ represent the latent utility of a vaccine profile with delay 
$t$ (months). The exponential discount rate $\delta$ is recovered from the relation 
$U(t) \propto e^{-\delta t}$. The hyperbolic parameter $\kappa$ is recovered from 
$U(t) \propto (1 + \kappa t)^{-1}$.

\subsection*{C.2 Estimation Results}

Based on the latest maximum-likelihood estimation of the Choice experiment data, we 
recover the following structural parameters:
\begin{itemize}
    \item \textbf{Hyperbolic Parameter ($\kappa = 0.225$)}: This indicates significant 
    short-run impatience, with utility declining sharply over the first several months 
    of delay.
    \item \textbf{Exponential Discount Rate ($\delta = 0.160$)}: This reflects a 
    steady decay in the perceived benefit of protection as the waiting interval increases.
\end{itemize}

These values indicate that a one-month delay is associated with a 16--22\% reduction in 
perceived utility relative to immediate vaccination. This high sensitivity to delay 
underscores the substantial behavioral burden imposed by waiting while infectious risk 
persists.


\subsection*{C.3 Summary Table}

\begin{table}[H]
\centering
\caption{Recovered Structural Discounting Parameters}
\label{tab:discount_rates}
\begin{tabular}{lcc}
\toprule
\textbf{Parameter} & \textbf{Value} & \textbf{Interpretation} \\
\midrule
Hyperbolic ($\kappa$) & 0.225 & Short-run impatience toward delay \\
Exponential ($\delta$) & 0.160 & Constant decay rate of protection value \\
\bottomrule
\end{tabular}
\end{table}

\clearpage

\end{document}



